\documentclass[11pt]{article}
\usepackage{amsmath,amssymb,amsthm}
\newtheorem{theorem}{Theorem}
\begin{document}
\title{When compact pole permutations have the same sign family}
\author{}\date{}\maketitle

\paragraph{Status and scope.}
This is a complete written audit of a proposed generic-query construction.
It is not a lower bound for the original HeadComplexity model and does not
settle either accepted asymptotic target. In particular, the query offsets
in this note are fixed while the pole scale tends to zero. Degenerate or
scale-dependent offsets, adaptive residue bases, and a construction at a
specified positive compact scale are outside its conclusion. The exact
companion checker tests a small instance; no Lean theorem is asserted.

Fix positive offsets $D_{hj}$, pole labels $p_1,\ldots,p_H$, and a
permutation $\pi$. Consider the linear score family on $q$ queries
\[
 F_j=c+\sum_{h=1}^H\frac{a_h}{D_{hj}-\epsilon p_{\pi(h)}}.
 \tag{1}
\]
Write $d=H+1$ and let $M_\pi(\epsilon)$ be its $q$ by $d$
coefficient matrix, including the constant column. Its common limit is
\[
 M_0(j,:)=(1,D_{1j}^{-1},\ldots,D_{Hj}^{-1}).
 \tag{2}
\]

\begin{theorem}[Uniform-minor stability]
Suppose $q\geq d$ and every $d$ by $d$ row minor of $M_0$ is nonzero.
There is $\epsilon_0>0$ such that, for every
$0<\epsilon<\epsilon_0$ and every pole permutation, the strict sign
vectors realized by (1) are exactly those realized by $M_0$.
Their number is
\[
 2\sum_{k=0}^{d-1}\binom{q-1}{k}.                 \tag{3}
\]
Consequently, taking the union over all $H!$ permutations produces no
additional strict sign vector in this interval.
\end{theorem}
\begin{proof}
First restrict $\epsilon$ so that all displayed denominators stay
positive. Each maximal minor is then continuous. There are finitely
many minors and permutations, and all their limits are nonzero.
Shrink the interval so that every minor retains its limiting sign.

Here is why these signs determine strict feasibility. Given a desired
sign vector $s$, Gordan's alternative says that
$\operatorname{diag}(s)Mv>0$ is infeasible exactly when there is a
nonzero $\lambda\geq0$ with
$\lambda^{\mathsf T}\operatorname{diag}(s)M=0$.
Choose a certificate of minimal support. Its support is a circuit:
otherwise a nonproportional dependence can be added with a small real
coefficient until one positive entry becomes zero, contradicting
minimality. Every $d$ rows are independent, so every circuit has
exactly $d+1$ rows. On an ordered support $i_0,\ldots,i_d$, its
dependence coefficients, up to a common nonzero factor, are
\[
 (-1)^k\det M_{\{i_0,\ldots,i_d\}\setminus\{i_k\}}.
\]
Whether multiplication by the desired signs makes these coefficients
all positive depends only on the maximal-minor signs. Thus strict
feasibility is unchanged. When $q=d$ the matrix is invertible and all
sign vectors are feasible, without any circuit test.

For completeness, the number in (3) is the number of chambers of the
$q$ central hyperplanes defined by the rows in $\mathbb R^d$.
The nonzero-minor assumption puts them in general position. Adding a
hyperplane creates as many chambers as the previous hyperplanes induce
on it, giving the recurrence
$R(q,d)=R(q-1,d)+R(q-1,d-1)$, with
$R(q,1)=2$ and $R(q,d)=2^q$ for $q\leq d$.
Pascal's identity gives (3). A strict sign cone is convex, so it cannot
contain two different chambers with the same sign vector.
\end{proof}

\paragraph{A fully explicit generic baseline.}
Choose $0<t_1<\cdots<t_q<1$ and
$D_{hj}=1/(1+t_j^h)$. Subtracting the constant column from the other
columns makes every maximal minor a Vandermonde determinant. Thus
the theorem applies to this positive, bounded family of offsets.
For $t_j=j/(q+1)$, $P=\max(1,\max_h|p_h|)$, and
$d=H+1$, one sufficient interval is
\[
 0<\epsilon<\min\left\{\frac1{8P},
 \frac{(q+1)^{-H(H+1)/2}}{16P\,d!\,d\,3^{d-1}}\right\}.
 \tag{4}
\]
Indeed the absolute value of every baseline minor is at least
$(q+1)^{-H(H+1)/2}$. Since $D_{hj}>1/2$, every perturbed reciprocal
differs from its baseline by at most $8\epsilon P$. If the maximum
entry error is $E\leq1$, expanding determinants as sums of products
and telescoping each product bounds their difference by
$d!dE3^{d-1}$. Formula (4) makes this less than half the lower bound.
The bias column has no error, so including it in this estimate only
enlarges the bound.

\paragraph{Asymptotic implication for this attempted construction.}
For $q$ of order $H\log H$, the fraction in (3), divided by $2^q$,
tends to zero. For example the elementary bound
$\sum_{k=0}^{d-1}\binom{q-1}{k}\leq
(e(q-1)/(d-1))^{d-1}$ applies when $1\leq d-1\leq q-1$.
The logarithm of the fraction is at most
$-q\log2+O(H\log\log H)$.
Thus a generic fixed-offset, arbitrarily-compact branch argument does
not provide the uniformly dense decoder needed by the checked batching
theorem. The very small interval (4) must not be confused with a bound
at every full-cube-admissible positive scale. Structured singular
limits, including the butterfly constructions investigated separately,
are not covered by this audit.
\end{document}
